\section{Software Tools and Libraries Used}

The implementation of the Medicare Bot system utilizes a combination of software libraries and frameworks for computer vision, backend development, and frontend visualization.

\begin{itemize}
\item \textbf{OpenCV}: Used for image processing and face detection.
\item \textbf{InsightFace}: Used for deep learning-based face recognition.
\item \textbf{ONNX Runtime}: Used for efficient execution of deep learning models.
\item \textbf{FastAPI}: Used to develop backend APIs for communication and data management.
\item \textbf{SQLite}: Used as the database for storing patients, medications, and logs.
\item \textbf{React}: Used to develop the frontend dashboard interface.
\item \textbf{Python}: Core programming language used for system implementation.
\end{itemize}

\section{Serial Communication Commands}

The Raspberry Pi communicates with the Arduino using serial commands to control the dispensing mechanism.

\begin{itemize}
\item \texttt{d1 -> Activate servo 1}
\item \texttt{d2 -> Activate servo 2}
\item \texttt{d3 -> Activate servo 3}
\end{itemize}

These commands are transmitted via USB serial communication and interpreted by the Arduino to perform the corresponding dispensing action.

\section{Sample API Request (Log Transmission)}

The Raspberry Pi sends log data to the backend server using an HTTP POST request with multipart/form-data.

\begin{verbatim}
POST /api/logs
Content-Type: multipart/form-data

patient_id=1
status=DISPENSED
compartment=2
confidence=0.87
image=file.jpg
\end{verbatim}

This request ensures both textual data and image files are transmitted correctly to the backend server.

\section{Configuration Parameters}

The system uses several configuration parameters to control performance and behavior:

\begin{itemize}
\item \texttt{RECOGNITION\_THRESHOLD = 0.45}
\item \texttt{UNCERTAIN\_THRESHOLD = 0.35}
\item \texttt{PROCESS\_EVERY\_N = 5}
\item \texttt{RESIZE\_SCALE = 0.5}
\item \texttt{CONFIRM\_SECONDS = 2.5}
\item \texttt{DISPENSE\_COOLDOWN\_SECONDS = 60}
\end{itemize}

These parameters are used to balance recognition accuracy, computational efficiency, and system safety.

\section{Software Project Structure}

The software system is organized into modular components as shown below:

\begin{verbatim}
Face_recognition/
|
+-- main.py
+-- database/
|   +-- db.py
+-- face/
|   +-- register.py
|   +-- recognize.py
|   +-- utils.py
+-- medication/
|   +-- scheduler.py
|   +-- dispenser.py
+-- tools/
|   +-- logger.py
+-- patients.db
\end{verbatim}

This modular structure improves maintainability, scalability, and readability of the system.

\section{System Design and Implementation Gallery}

\subsection{3D Design Models}

\begin{figure}[H]
    \centering
    % First Row: Images 16 and 17
    \includegraphics[width=0.35\linewidth]{16.png}
    \hspace{0.5cm} % Adds horizontal space between images
    \includegraphics[width=0.35\linewidth]{17.png} 
    \\[10pt] % Breaks the line and adds 10pt vertical spacing

    % Second Row: Images 18 and 19
    \includegraphics[width=0.35\linewidth]{18.png}
    \hspace{0.5cm}
    \includegraphics[width=0.35\linewidth]{19.png}
    \\[10pt]

    % Third Row: Image 20 (Centered)
    \includegraphics[width=0.35\linewidth]{20.png}

    \caption{3D designs}
    \label{fig:3d_designs}
\end{figure}

\subsection{Code Overview}

\begin{figure}[H]
    \centering
    % First Row: Images 16 and 17
    \includegraphics[width=0.35\linewidth]{code1.png}
    \hspace{0.5cm} % Adds horizontal space between images
    \includegraphics[width=0.35\linewidth]{code2.png} 
    \\[10pt] % Breaks the line and adds 10pt vertical spacing

    % Second Row: Images 18 and 19
    \includegraphics[width=0.2\linewidth]{code3.png}
    \hspace{0.5cm}
    \includegraphics[width=0.35\linewidth]{code4.png}
    \\[10pt]

    % Third Row: Image 20 (Centered)
    \includegraphics[width=0.35\linewidth]{codex.png}

    \caption{Code Overview}
    \label{fig:Code Overview}
\end{figure}
\subsection{Pi Configuration}

\begin{figure}[H]
 \centering
    % First Row: Images 16 and 17
    \includegraphics[width=0.35\linewidth]{pi.jpg}
    \hspace{0.5cm} % Adds horizontal space between images
    \includegraphics[width=0.35\linewidth]{pi1.jpg} 
    \\[10pt] % Breaks the line and adds 10pt vertical spacing

    % Second Row: Images 18 and 19
    \includegraphics[width=0.35\linewidth]{pi2.jpg}
      \includegraphics[width=0.35\linewidth]{pi3.jpg}
      \includegraphics[width=0.35\linewidth]{pi4.jpg}
      \includegraphics[width=0.35\linewidth]{pi5.jpg}
      \includegraphics[width=0.35\linewidth]{pi6.jpg}
      \includegraphics[width=0.35\linewidth]{pi7.jpg}
                
    \caption{Pi Configuration}
    \label{fig:Pi Configuration}
\end{figure}

\subsection{Final System}

\begin{figure}[H]
 \centering
    % First Row: Images 16 and 17
    \includegraphics[width=0.35\linewidth]{bot.jpg}
    \hspace{0.5cm} % Adds horizontal space between images
    \includegraphics[width=0.35\linewidth]{bot2.jpg} 
    \\[10pt] % Breaks the line and adds 10pt vertical spacing

    % Second Row: Images 18 and 19
    \includegraphics[width=0.35\linewidth]{bot3.jpg}
    \caption{Final Medicare Bot System Implementation}
    \label{fig:Final Medicare Bot System Implementation}
\end{figure}